\textbf{Phần 1. Trắc nghiệm lựa chọn}
\setcounterpageref{ex}{0}

\begin{ex} %[5X11V67]
Biểu đồ quạt tròn biểu diễn kết quả điều tra về sở thích thể thao của $300$ học sinh khối 5: Môn Bơi lội chiếm $35\%$, môn Bóng đá chiếm $40\%$, môn Cầu lông chiếm $15\%$, môn Cờ vua chiếm $10\%$. Số học sinh thích môn Bóng đá nhiều hơn số học sinh thích môn Bơi lội là
\choice
{$12\text{ bạn}$}
{\True $15\text{ bạn}$}
{$18\text{ bạn}$}
{$20\text{ bạn}$}
\loigiai{
Hiệu tỉ số phần trăm giữa môn Bóng đá và Bơi lội là:
$$40\% - 35\% = 5\%.$$
Số học sinh chênh lệch là:
$$300 \times 5\% = 15\text{ (bạn)}.$$
\textbf{Đáp án của câu hỏi là B}
}
\end{ex}

\begin{ex} %[5X11V67]
Gieo một con xúc xắc $100$ lần, mặt có số chấm chia hết cho $2$ xuất hiện $54$ lần. Tỉ số phần trăm số lần xuất hiện mặt có số chấm lẻ là
\choice
{$44\%$}
{\True $46\%$}
{$48\%$}
{$52\%$}
\loigiai{
Số lần xuất hiện mặt lẻ là:
$$100 - 54 = 46\text{ (lần)}.$$
Tỉ số phần trăm là:
$$46 : 100 \times 100\% = 46\%.$$
\textbf{Đáp án của câu hỏi là B}
}
\end{ex}

\begin{ex} %[5X11V67]
Một gia đình ghi lại số tiền chi cho ăn uống trong $4$ tuần: Tuần 1 chi $1\,400\,000\text{ đồng}$, Tuần 2 chi $1\,600\,000\text{ đồng}$, Tuần 3 chi $1\,500\,000\text{ đồng}$, Tuần 4 chi $1\,900\,000\text{ đồng}$. Trung bình mỗi tuần gia đình đó chi bao nhiêu tiền cho ăn uống?
\choice
{$1\,550\,000\text{ đồng}$}
{\True $1\,600\,000\text{ đồng}$}
{$1\,650\,000\text{ đồng}$}
{$1\,700\,000\text{ đồng}$}
\loigiai{
Tổng số tiền chi trong $4$ tuần là:
$$1\,400\,000 + 1\,600\,000 + 1\,500\,000 + 1\,900\,000 = 6\,400\,000\text{ (đồng)}.$$
Trung bình mỗi tuần chi là:
$$6\,400\,000 : 4 = 1\,600\,000\text{ (đồng)}.$$
\textbf{Đáp án của câu hỏi là B}
}
\end{ex}

\begin{ex} %[5X11V67]
Biểu đồ quạt tròn biểu thị tỉ lệ tiêu thụ các loại hoa quả tại một siêu thị trong ngày: Táo ($30\%$), Nho ($25\%$), Dưa hấu ($35\%$), Cam ($10\%$). Biết số lượng dưa hấu bán được là $140\text{ kg}$. Tổng số ki-lô-gam của cả $4$ loại hoa quả bán được là
\choice
{$350\text{ kg}$}
{\True $400\text{ kg}$}
{$420\text{ kg}$}
{$450\text{ kg}$}
\loigiai{
Tổng khối lượng hoa quả bán được là:
$$140 : 35\% = 400\text{ (kg)}.$$
\textbf{Đáp án của câu hỏi là B}
}
\end{ex}

\begin{ex} %[5X11V67]
Tung hai đồng xu cùng một lúc $60$ lần, thấy có $33$ lần xuất hiện hai mặt giống nhau (cùng Sấp hoặc cùng Ngửa). Tỉ số số lần xuất hiện hai mặt khác nhau so với tổng số lần tung viết dưới dạng phân số tối giản là
\choice
{$\dfrac{11}{20}$}
{\True $\dfrac{9}{20}$}
{$\dfrac{7}{20}$}
{$\dfrac{3}{10}$}
\loigiai{
Số lần xuất hiện hai mặt khác nhau là:
$$60 - 33 = 27\text{ (lần)}.$$
Tỉ số là:
$$\dfrac{27}{60} = \dfrac{9}{20}.$$
\textbf{Đáp án của câu hỏi là B}
}
\end{ex}

\begin{ex} %[5X11V67]
Cho bảng số liệu về số cây trồng được của ba tổ: Tổ 1 trồng $42$ cây, Tổ 2 trồng $48$ cây. Tổ 3 trồng được số cây bằng trung bình cộng số cây của cả $3$ tổ. Số cây Tổ 3 trồng được là
\choice
{$44\text{ cây}$}
{\True $45\text{ cây}$}
{$46\text{ cây}$}
{$48\text{ cây}$}
\loigiai{
Vì Tổ 3 trồng bằng trung bình cộng cả 3 tổ nên Tổ 3 trồng đúng bằng trung bình cộng của Tổ 1 và Tổ 2:\\
$$(42 + 48) : 2 = 45\text{ (cây)}.$$
\textbf{Đáp án của câu hỏi là B}
}
\end{ex}

\begin{ex} %[5X11V67]
Biểu đồ quạt tròn cho biết cơ cấu nguồn vốn một dự án nhỏ: Vốn tự có chiếm $45\%$, Vốn vay ngân hàng chiếm $35\%$, Vốn kêu gọi đầu tư chiếm $20\%$. Biết số vốn kêu gọi đầu tư là $60\,000\,000\text{ đồng}$. Số vốn tự có của dự án đó là
\choice
{$120\,000\,000\text{ đồng}$}
{\True $135\,000\,000\text{ đồng}$}
{$150\,000\,000\text{ đồng}$}
{$140\,000\,000\text{ đồng}$}
\loigiai{
Tổng nguồn vốn của dự án là:
$$60\,000\,000 : 20\% = 300\,000\,000\text{ (đồng)}.$$
Số vốn tự có là:
$$300\,000\,000 \times 45\% = 135\,000\,000\text{ (đồng)}.$$
\textbf{Đáp án của câu hỏi là B}
}
\end{ex}

\begin{ex} %[5X11V67]
Một hộp kín đựng $4$ loại kẹo có số lượng bằng nhau: Kẹo dâu, Kẹo nho, Kẹo cam, Kẹo bạc hà. Bạn Minh bốc ngẫu nhiên $1$ chiếc kẹo, ghi lại loại rồi bỏ lại vào hộp. Bạn thực hiện $80$ lần và có $28$ lần bốc được kẹo dâu. Tỉ số phần trăm thực nghiệm bốc được kẹo dâu cao hơn so với tỉ lệ lí thuyết ($25\%$) là
\choice
{$8\%$}
{\True $10\%$}
{$12\%$}
{$15\%$}
\loigiai{
Tỉ số phần trăm thực nghiệm bốc được kẹo dâu là:
$$28 : 80 \times 100\% = 35\%.$$
Cao hơn tỉ lệ lí thuyết số phần trăm là:
$$35\% - 25\% = 10\%.$$
\textbf{Đáp án của câu hỏi là B}
}
\end{ex}

\begin{ex} %[5X11V67]
Thống kê điểm số môn Tiếng Anh của $30$ học sinh: Có $6$ bạn đạt điểm 10, $9$ bạn đạt điểm 9, $12$ bạn đạt điểm 8, số bạn còn lại đạt điểm 7. Tỉ số phần trăm các bạn đạt điểm từ 9 trở lên là
\choice
{$45\%$}
{\True $50\%$}
{$55\%$}
{$60\%$}
\loigiai{
Số bạn đạt điểm 9 và 10 là:
$$6 + 9 = 15\text{ (bạn)}.$$
Tỉ số phần trăm là:
$$15 : 30 \times 100\% = 50\%.$$
\textbf{Đáp án của câu hỏi là B}
}
\end{ex}

\begin{ex} %[5X11V67]
Biểu đồ quạt tròn biểu thị diện tích các loại cây ăn trái của một nhà vườn: Vải thiều $40\%$, Nhãn lồng $30\%$, Mãng cầu $20\%$, Thanh long $10\%$. Diện tích trồng Vải thiều gấp mấy lần diện tích trồng Thanh long?
\choice
{$3\text{ lần}$}
{\True $4\text{ lần}$}
{$2\text{ lần}$}
{$5\text{ lần}$}
\loigiai{
Tỉ số giữa tỉ lệ diện tích Vải thiều và Thanh long là:
$$40\% : 10\% = 4\text{ (lần)}.$$
\textbf{Đáp án của câu hỏi là B}
}
\end{ex}

\begin{ex} %[5X11V67]
Một xạ thủ bắn $150$ phát đạn, có $126$ phát bắn trúng vòng $9$ và vòng $10$. Tỉ số phần trăm số phát bắn không đạt điểm $9$ hoặc $10$ là
\choice
{$14\%$}
{\True $16\%$}
{$18\%$}
{$20\%$}
\loigiai{
Số phát bắn không đạt vòng 9 hoặc 10 là:
$$150 - 126 = 24\text{ (phát)}.$$
Tỉ số phần trăm là:
$$24 : 150 \times 100\% = 16\%.$$
\textbf{Đáp án của câu hỏi là B}
}
\end{ex}

\begin{ex} %[5X11V67]
Lớp 5B có $40$ học sinh. Biểu đồ quạt tròn cho biết có $65\%$ số bạn tham gia đội nghi thức Đội, $20\%$ tham gia đội văn nghệ xung kích, số còn lại tham gia đội cờ đỏ. Số học sinh tham gia đội cờ đỏ là
\choice
{$5\text{ bạn}$}
{\True $6\text{ bạn}$}
{$7\text{ bạn}$}
{$8\text{ bạn}$}
\loigiai{
Tỉ số phần trăm học sinh tham gia đội cờ đỏ là:
$$100\% - (65\% + 20\%) = 15\%.$$
Số học sinh tham gia đội cờ đỏ là:
$$40 \times 15\% = 6\text{ (bạn)}.$$
\textbf{Đáp án của câu hỏi là B}
}
\end{ex}

\textbf{Phần 2. Trắc nghiệm đúng sai}
\setcounterpageref{ex}{0}

\begin{ex} %[5X11V67]
Biểu đồ quạt tròn biểu diễn kết quả đánh giá sự hài lòng của $400$ khách hàng đối với một sản phẩm: Rất hài lòng ($50\%$), Hài lòng ($35\%$), Bình thường ($10\%$), Không hài lòng ($5\%$). Xét tính đúng/sai của các nhận định sau:
\choiceTF
{\True Có $200$ khách hàng đánh giá Rất hài lòng}
{\True Số khách hàng đánh giá Hài lòng là $140$ người}
{Số khách hàng đánh giá Không hài lòng là $15$ người}
{\True Tổng số khách hàng đánh giá từ Hài lòng trở lên chiếm $85\%$}
\loigiai{
\begin{itemchoice}
\item Đúng. $400 \times 50\% = 200$ khách hàng.
\item Đúng. $400 \times 35\% = 140$ khách hàng.
\item Sai. Số khách Không hài lòng là $400 \times 5\% = 20$ người (không phải 15 người).
\item Đúng. Từ Hài lòng trở lên: $50\% + 35\% = 85\%$.
\end{itemchoice}
}
\end{ex}

\begin{ex} %[5X11V67]
Một bạn học sinh quay đĩa số $50$ lần và ghi nhận số lần xuất hiện các màu trên đĩa: Màu Đỏ ($15$ lần), Màu Xanh ($20$ lần), Màu Vàng ($10$ lần), Màu Tím ($5$ lần). Xét tính đúng/sai của các nhận định sau:
\choiceTF
{\True Màu Xanh xuất hiện nhiều lần nhất với tỉ số phần trăm là $40\%$}
{\True Màu Tím xuất hiện ít lần nhất với tỉ số phần trăm là $10\%$}
{Tỉ số phần trăm số lần xuất hiện Màu Đỏ là $25\%$}
{\True Tổng số lần xuất hiện Màu Đỏ và Màu Vàng bằng một nửa tổng số lần quay}
\loigiai{
\begin{itemchoice}
\item Đúng. $20 : 50 = 40\%$ (nhiều nhất).
\item Đúng. $5 : 50 = 10\%$ (ít nhất).
\item Sai. Màu Đỏ chiếm $15 : 50 = 30\%$ (không phải 25\%).
\item Đúng. Đỏ + Vàng $= 15 + 10 = 25$ lần $= \dfrac{25}{50} = \dfrac{1}{2}$ (một nửa).
\end{itemchoice}
}
\end{ex}

\textbf{Phần 3. Câu trả lời ngắn}
\setcounterpageref{ex}{0}

\begin{ex} %[5X11V67]
Một trường tiểu học có $800$ học sinh. Biểu đồ quạt tròn cho biết có $45\%$ số học sinh đạt danh hiệu Học sinh Xuất sắc. Hỏi trường đó có bao nhiêu bạn đạt danh hiệu Học sinh Xuất sắc?
\shortans{360}
\loigiai{
Số học sinh đạt danh hiệu Học sinh Xuất sắc là:
$$800 \times 45\% = 360\text{ (học sinh)}.$$
}
\end{ex}

\begin{ex} %[5X11V67]
Trong một đợt thi bắn bia, xạ thủ bắn $60$ phát và đạt tỉ lệ trúng bia là $95\%$. Hỏi xạ thủ đó bắn trúng bia bao nhiêu phát?
\shortans{57}
\loigiai{
Số phát bắn trúng bia là:
$$60 \times 95\% = 57\text{ (phát)}.$$
}
\end{ex}

\begin{ex} %[5X11V67]
Gieo một con xúc xắc $120$ lần, sự kiện “xuất hiện mặt có số chấm là hợp số” ($4$ hoặc $6$ chấm) xuất hiện $38$ lần. Hỏi số lần xuất hiện các mặt không phải hợp số ($1, 2, 3, 5$) là bao nhiêu lần?
\shortans{82}
\loigiai{
Số lần xuất hiện các mặt không phải hợp số là:
$$120 - 38 = 82\text{ (lần)}.$$
}
\end{ex}

\begin{ex} %[5X11V67]
Biểu đồ quạt tròn biểu thị tỉ lệ phương tiện của cư dân một khu phố: Xe máy chiếm $70\%$, Xe đạp chiếm $12\%$, Ô tô chiếm $18\%$. Biết khu phố có $54$ chiếc ô tô. Hỏi khu phố đó có tất cả bao nhiêu chiếc xe máy?
\shortans{210}
\loigiai{
Tổng số phương tiện của khu phố là:
$$54 : 18\% = 300\text{ (chiếc)}.$$
Số chiếc xe máy của khu phố là:
$$300 \times 70\% = 210\text{ (chiếc)}.$$
}
\end{ex}

\begin{ex} %[5X11V67]
Một cửa hàng nhập về $500\text{ kg}$ hoa quả gồm $3$ loại: Xoài, Cam và Táo. Biểu đồ quạt tròn cho biết: Xoài chiếm $36\%$, Cam chiếm $44\%$, còn lại là Táo. Hỏi cửa hàng nhập về bao nhiêu ki-lô-gam táo?
\shortans{100}
\loigiai{
Tỉ số phần trăm khối lượng táo là:
$$100\% - (36\% + 44\%) = 20\%.$$
Khối lượng táo nhập về là:
$$500 \times 20\% = 100\text{ (kg)}.$$
}
\end{ex}

\begin{ex} %[5X11V67]
Điểm kiểm tra của $4$ bạn học sinh là: $8, 9, 9, 10$. Hỏi điểm trung bình cộng của $4$ bạn đó là bao nhiêu? (Nếu là số thập phân thì nhân $10$, ví dụ $8{,}5$ thì nhập $85$, hoặc tính chẵn: $8+9+9+10 = 36 \Rightarrow 36 : 4 = 9$).
\shortans{9}
\loigiai{
Tổng số điểm của 4 bạn là:
$$8 + 9 + 9 + 10 = 36.$$
Điểm trung bình cộng là:
$$36 : 4 = 9.$$
}
\end{ex}
